\label{chap:from-themes-to-requirements}

This chapter turns the themes of Chapter~\ref{chap:formative-study} into features and testable requirements. It follows the logic of human-centred design, in which design rests on an explicit understanding of users, tasks and environments and proceeds through specifying the user requirements, producing design solutions, and evaluating them against those requirements~\cite{iso9241_210}, and that of requirements engineering, which identifies stakeholders and their needs and documents them in a form amenable to analysis, communication and subsequent implementation~\cite[p.~37]{nuseibeh_requirements_2000}.

Together with Chapter~\ref{chap:formative-study}, this chapter answers the first research question:

\begin{quote}
\textit{What are students' needs and resulting requirements for a hybrid graph-table interface that supports programme-level study planning and understanding of curriculum structure and course dependencies?}
\end{quote}

Its outcome is the features-and-requirements specification that the Design and Implementation chapters build (Chapter~\ref{chap:design}, Chapter~\ref{chap:implementation}) and the evaluation reports against (Chapter~\ref{chap:evaluation}).

\section{From Themes to Requirements}
\label{sec:derivation-process}

The starting point was the final theme set of Chapter~\ref{chap:formative-study}: 64 themes, each a design-relevant need category grounded in collated interview excerpts and linked throughout to that evidence, so that a requirement can be followed backwards to its qualitative origin.

Each theme was first read, together with the initial excerpts constituting it, to formulate a design implication, the guiding question being what the system must support if the need is to be addressed. Each theme, rather than a single utterance, is the operative expression of a stakeholder need: elicited material has to be interpreted, analysed and modelled before it can stand as a requirement~\cite[p.~39]{nuseibeh_requirements_2000}, and the theme is the unit at which that interpretation was recorded, with its excerpts behind it. Each theme then yielded one or more candidate \gterm{requirement}{requirements}, written as statements of what the system must support rather than how it should be built. Related requirements were grouped into \emph{\gterm{feature}{features}}, a feature being one coherent capability of the product. Where several themes converged on one capability they shared a feature; where one theme implied different capabilities it fed more than one, so a requirement count is not tied one-to-one to its theme. Where something was deliberately left out of a feature, a short delineation records it, consistent with requirements-engineering practice, in which not all elicited needs are taken forward unchanged but must be bounded in a form manageable for development.

The chain this produces is \emph{interview evidence $\rightarrow$ code $\rightarrow$ theme $\rightarrow$ requirement $\rightarrow$ feature}, matching the derivation order above and the codebook Chapter~\ref{chap:formative-study} produced (Section~\ref{sec:output-traceability}). Linking a requirement back to what motivated it before any specification existed is what Gotel and Finkelstein call pre-requirements-specification traceability~\cite[p.~96]{gotel_analysis_1994}. Each feature accordingly names the themes it rests on by identifier, and each theme's excerpts are reproduced in Section~\ref{app:codebook}, so the chain can be followed in either direction.

\section{Resulting Features}
\label{sec:resulting-features}

The theme set yielded 13 features carrying 55 numbered requirements between them. The five scope boundaries OOS1 to OOS5 defined in the Methodology chapter (Section~\ref{sec:overall-scope}) apply throughout, and are referred to by those numbers here and in later chapters. Each feature carries a total importance score (Section~\ref{sec:prevalence-ranking}), the summed prevalence of its associated themes, which was used to prioritise early implementation attention. Table~\ref{tab:feature-summary} lists the result in descending order of that score. The full specifications are given in Section~\ref{app:feature-specs}, each in the same structure: title and identifier, importance score, associated themes, description and user evidence, derived requirements, and a delineation where one applies; they are cited from the Design chapter onward as, e.g., Feature~\ref{feat:004}, Req.~1.

The score is a sum, not a mean, so it rewards breadth: a feature seven themes converge on outranks one that two themes raise more insistently. Normalising by theme count would order the specification differently. The order was used once, to sequence implementation work, and no finding in Chapter~\ref{chap:evaluation} or Chapter~\ref{chap:discussion} depends on it.

The 13 features fall into three groups. Four are the planning core: getting curriculum data into the system without manual transcription (FEATURE-001), showing progress against the programme at three granularities (FEATURE-002), enforcing the curriculum's own rules (FEATURE-003), and moving a course from exploration in the graph to commitment in the table (FEATURE-004). Together they define the hybrid graph--table concept the rest of the thesis builds and tests. Four more add the judgement layer participants described around that core: semester feasibility (FEATURE-005), explainable recommendation (FEATURE-007), stage-appropriate onboarding (FEATURE-006), and replanning under disruption (FEATURE-008). The remaining five concern how the representation behaves once it is populated: density control (FEATURE-009), graph interpretability (FEATURE-010), visual and affective load (FEATURE-011), non-destructive recommendation display (FEATURE-012), and user-defined spatial grouping (FEATURE-013). The three groups run in descending score order, which is why implementation began with the first.

Two features contain requirements that were derived from the data but not carried into the implemented system, because of resource and time constraints during development: FEATURE-005, Req.~3 (self-assessed readiness profiles) and FEATURE-013, Reqs.~1, 3, and~4 (named, cross-view groups). Both are stated in the respective Delineation in Section~\ref{app:feature-specs} and are taken up as shortfalls in Chapter~\ref{chap:discussion}. Every other exclusion recorded in a Delineation follows from a scope boundary, OOS2 to OOS4, or from a design decision stated there, not from what was built.

\begin{table}[htpb]
    \centering
    \footnotesize
    \begin{tabular}{@{}l>{\raggedright\arraybackslash}p{4.5cm}c>{\raggedright\arraybackslash}p{3.3cm}c@{}}
        \toprule
        \textbf{ID} & \textbf{Feature} & \textbf{Score} & \textbf{Associated themes} & \textbf{Reqs.} \\
        \midrule
        FEATURE-001 & Unified Data \& Tool Integration & 20 & T9, T10, T11, T12, T15, T37, T62 & 5 \\
        FEATURE-002 & Multi-Level Roadmap \& Progress \gterm{dashboard}{Dashboard} & 19 & T1, T16, T18, T19, T57 & 6 \\
        FEATURE-003 & Curriculum Rules \& Compliance Engine & 18 & T13, T17, T30, T31, T33, T63 & 5 \\
        FEATURE-004 & Graph-to-Table Planning Workflow & 15 & T20, T21, T44, T47, T61 & 6 \\
        FEATURE-005 & Feasibility Planning & 14 & T6, T35, T48, T49 & 3 \\
        FEATURE-007 & Evidence-Based, Explainable Recommendation Support & 13 & T23, T24, T45, T52 & 6 \\
        FEATURE-006 & Onboarding \& Guided Planning Transition & 12 & T5, T7, T59 & 2 \\
        FEATURE-008 & Constraint-Aware Replanning \& Disruption Handling & 9 & T4, T46, T56 & 4 \\
        FEATURE-009 & Overwhelm-Safe Browsing & 6 & T42, T50, T64 & 3 \\
        FEATURE-010 & Graph Interpretability & 5 & T27, T32, T40, T41 & 4 \\
        FEATURE-011 & Motivating Visual UI \& Emotional Load Reduction & 5 & T8, T54, T58 & 4 \\
        FEATURE-012 & Recommendation Presentation UX & 4 & T43, T55 & 3 \\
        FEATURE-013 & Spatial Organisation \& Grouping Controls & 3 & T28, T29 & 4 \\
        \bottomrule
    \end{tabular}
    \caption{The 13 features, in descending order of importance score, with the themes each was derived from and the number of requirements each carries. The full specifications are in Section~\ref{app:feature-specs}.}
    \label{tab:feature-summary}
\end{table}
